\chapter{Vectors and Lines in the Plane}\label{ch:g11:vect}

\omterm{def:g11:vect:vector}{Vectors} encode displacements; lines are their trajectories. The key new
tool of this chapter is the \emph{\omterm{def:g11:vect:det}{determinant}}, a single number that
decides whether two \omterm{def:g11:vect:vector}{vectors} are \omterm{def:g11:vect:collinear}{collinear} --- and, from there, produces
\omterm{def:g10:algebra:equation}{equations} of lines and settles alignment and parallelism questions by pure
computation. The same ideas extend to space in \cref{ch:g12:space}.

\section{Vectors: a recap}

\begin{definition}[Vector, coordinates]\label{def:g11:vect:vector}
A \emph{vector}\index{vector} $\vec u$ represents a displacement: all
arrows with the same direction, orientation and length represent the same
vector. In a \omterm{def:g10:coordgeom:system}{coordinate system}, $\vec u\,(x, y)$ records the displacement
($x$ horizontally, $y$ vertically); for points $A(x_A, y_A)$ and
$B(x_B, y_B)$,
\[
\vect{AB}\,\bigl(x_B - x_A,\ y_B - y_A\bigr).
\]
Vectors add coordinate-wise, and $\lambda\vec u$ has \omterm{def:g10:coordgeom:system}{coordinates}
$(\lambda x, \lambda y)$.
\end{definition}

\begin{proposition}[Midpoint]\label{prop:g11:vect:midpoint}
The \omterm{prop:g10:coordgeom:midpoint}{midpoint} $M$ of the segment $[AB]$ has \omterm{def:g10:coordgeom:system}{coordinates}
$\left(\frac{x_A + x_B}{2},\ \frac{y_A + y_B}{2}\right)$.
\end{proposition}

\begin{proof}
$M$ is characterized by $\vect{AM} = \frac12\vect{AB}$: coordinate-wise,
$x_M - x_A = \frac12(x_B - x_A)$, so $x_M = \frac{x_A + x_B}{2}$, and
similarly for $y_M$.
\end{proof}

\section{Collinearity and the determinant}

\begin{definition}[Collinear vectors]\label{def:g11:vect:collinear}
Two \omterm{def:g11:vect:vector}{vectors} are \emph{collinear}\index{collinear vectors} if one is a
scalar multiple of the other: $\vec v = \lambda \vec u$ for some
$\lambda \in \R$ (or $\vec u = \vec 0$). Geometrically: they have the same
direction, or one of them is zero.
\end{definition}

\begin{definition}[Determinant]\label{def:g11:vect:det}
The \emph{determinant}\index{determinant} of $\vec u\,(x, y)$ and
$\vec v\,(x', y')$ is
\[
\det(\vec u, \vec v) =
\begin{vmatrix} x & x' \\ y & y' \end{vmatrix}
= x y' - x' y .
\]
\end{definition}

\begin{theorem}[Collinearity criterion]\label{thm:g11:vect:collinearity}
Two \omterm{def:g11:vect:vector}{vectors} $\vec u$ and $\vec v$ are \omterm{def:g11:vect:collinear}{collinear} if and only if
$\det(\vec u, \vec v) = 0$.
\end{theorem}

\begin{proof}
($\Rightarrow$) If $\vec v = \lambda\vec u$, then $x' = \lambda x$ and
$y' = \lambda y$, so
$\det(\vec u, \vec v) = x(\lambda y) - (\lambda x)y = 0$; if
$\vec u = \vec 0$, the \omterm{def:g11:vect:det}{determinant} is trivially $0$.

($\Leftarrow$) Suppose $xy' - x'y = 0$ and $\vec u \neq \vec 0$, say
$x \neq 0$ (the case $y \neq 0$ is symmetric). Set
$\lambda = \frac{x'}{x}$. Then $x' = \lambda x$ by construction, and the
relation $xy' = x'y = \lambda x y$ gives, after dividing by $x \neq 0$,
$y' = \lambda y$. Hence $\vec v = \lambda \vec u$.
\end{proof}

\begin{example}\label{ex:g11:vect:collinear}
$\vec u\,(3, -2)$ and $\vec v\,(-6, 4)$:
$\det = 3 \times 4 - (-6) \times (-2) = 12 - 12 = 0$, \omterm{def:g11:vect:collinear}{collinear} (indeed
$\vec v = -2\vec u$). $\vec u\,(3, -2)$ and $\vec w\,(2, 1)$:
$\det = 3 + 4 = 7 \neq 0$, not \omterm{def:g11:vect:collinear}{collinear}.
\end{example}

\begin{method}[Alignment of three points]\label{met:g11:vect:alignment}
$A$, $B$, $C$ are aligned if and only if $\vect{AB}$ and $\vect{AC}$ are
\omterm{def:g11:vect:collinear}{collinear}: compute both \omterm{def:g11:vect:vector}{vectors} and check
$\det\bigl(\vect{AB}, \vect{AC}\bigr) = 0$.
\end{method}

\section{Cartesian equations of lines}

\begin{definition}[Direction vector]\label{def:g11:vect:direction}
A \emph{direction vector}\index{direction vector} of a line $d$ is any
nonzero \omterm{def:g11:vect:vector}{vector} $\vec u$ such that $\vect{AB}$ is \omterm{def:g11:vect:collinear}{collinear} with $\vec u$
for all points $A, B$ of $d$.
\end{definition}

\begin{theorem}[Cartesian equation]\label{thm:g11:vect:cartesian}
A line with \omterm{def:g11:vect:direction}{direction vector} $\vec u\,(-b, a)$ admits an \omterm{def:g10:algebra:equation}{equation} of the
form
\[
ax + by + c = 0
\qquad (a, b) \neq (0, 0).
\]
Conversely, every such \omterm{def:g10:algebra:equation}{equation} describes a line with \omterm{def:g11:vect:direction}{direction vector}
$\vec u\,(-b, a)$.
\end{theorem}

\begin{proof}
Let $d$ pass through $A(x_A, y_A)$ with direction $\vec u\,(-b, a)$. A
point $M(x, y)$ lies on $d$ exactly when $\vect{AM}$ is \omterm{def:g11:vect:collinear}{collinear} with
$\vec u$, i.e.\ (\cref{thm:g11:vect:collinearity})
\[
0 = \det\bigl(\vect{AM}, \vec u\bigr)
= (x - x_A)\,a - (-b)(y - y_A)
= ax + by - (a x_A + b y_A),
\]
an \omterm{def:g10:algebra:equation}{equation} of the announced form with $c = -(a x_A + b y_A)$.
Conversely, the solutions of $ax + by + c = 0$ with, say, $b \neq 0$ are
the points $\left(x, -\frac{a x + c}{b}\right)$: subtracting two of them
shows every chord is \omterm{def:g11:vect:collinear}{collinear} with $(-b, a)$, and one solution always
exists --- it is the line through it directed by $(-b, a)$.
\end{proof}

\begin{method}[Line through two points]\label{met:g11:vect:twopoints}
To find an \omterm{def:g10:algebra:equation}{equation} of the line $(AB)$: compute the \omterm{def:g11:vect:direction}{direction vector}
$\vect{AB}\,(x_B - x_A, y_B - y_A)$; match it with $(-b, a)$, i.e.\ take
$a = y_B - y_A$ and $b = -(x_B - x_A)$; determine $c$ by plugging in the
\omterm{def:g10:coordgeom:system}{coordinates} of $A$.
\end{method}

\begin{example}\label{ex:g11:vect:twopoints}
Line through $A(1, 2)$ and $B(4, 3)$: $\vect{AB}\,(3, 1)$, so
$a = 1$, $b = -3$, and the \omterm{def:g10:algebra:equation}{equation} is $x - 3y + c = 0$. Plugging in $A$:
$1 - 6 + c = 0$, so $c = 5$. \omterm{def:g10:algebra:equation}{Equation}: $x - 3y + 5 = 0$. Check with $B$:
$4 - 9 + 5 = 0$.
\end{example}

\begin{remark}
When $b \neq 0$, the \omterm{def:g10:algebra:equation}{equation} can be solved for $y$: $y = mx + p$ with
\omterm{def:g10:reffunc:affine}{slope} $m = -\frac ab$. The cartesian form $ax + by + c = 0$ is more
general: it also covers vertical lines ($b = 0$), which have no \omterm{def:g10:reffunc:affine}{slope}.
\end{remark}

\begin{omfigure}
\begin{tikzpicture}
\begin{axis}[omaxis,
  width=8.6cm, height=5.8cm,
  xmin=-1.6, xmax=5.6, ymin=-0.9, ymax=4.4,
  xtick={1,4}, ytick={2,3}]
  \addplot[omDef, thick, domain=-1.4:5.4] {(x+5)/3};
  \fill (axis cs:1,2) circle (1.5pt)
    node[above left, font=\small] {$A$};
  \fill (axis cs:4,3) circle (1.5pt)
    node[above left, font=\small] {$B$};
  \draw[-Stealth, omThm, thick] (axis cs:1,2) -- (axis cs:2.5,2.5)
    node[midway, below right, font=\small] {$\vec u$};
\end{axis}
\end{tikzpicture}

{\small The line $x - 3y + 5 = 0$ of \cref{ex:g11:vect:twopoints}, its
\omterm{def:g11:vect:direction}{direction vector} $\vec u = \vect{AB}$ (red, here scaled by $\frac12$),
and the two points used to build the \omterm{def:g10:algebra:equation}{equation}.}
\end{omfigure}

\section{Parametric representation}

\begin{definition}[Parametric representation]\label{def:g11:vect:parametric}
The line through $A(x_A, y_A)$ with direction $\vec u\,(\alpha, \beta)$ is
the set of points
\[
M\bigl(x_A + t\alpha,\ y_A + t\beta\bigr), \qquad t \in \R .
\]
The number $t$ is the \emph{parameter}: each value of $t$ produces one
point of the line, as if travelling along it at constant speed $\vec u$.
\end{definition}

\begin{example}
The line of \cref{ex:g11:vect:twopoints} is also
$\{(1 + 3t,\ 2 + t) : t \in \R\}$: $t = 0$ gives $A$, $t = 1$ gives $B$.
Eliminating $t$ ($t = y - 2$, so $x = 1 + 3(y-2)$) recovers
$x - 3y + 5 = 0$.
\end{example}

\section{Relative positions of two lines}

\begin{proposition}[Parallel or intersecting]\label{prop:g11:vect:positions}
Two lines with \omterm{def:g11:vect:direction}{direction vectors} $\vec u$ and $\vec v$ are parallel if and
only if $\det(\vec u, \vec v) = 0$. In \omterm{def:g10:algebra:equation}{equation} form: $d : ax + by + c =
0$ and $d' : a'x + b'y + c' = 0$ are parallel if and only if
$ab' - a'b = 0$. Non-parallel lines meet in exactly one point.
\end{proposition}

\begin{proof}
Parallelism means \omterm{def:g11:vect:collinear}{collinear} directions, which is
\cref{thm:g11:vect:collinearity}; with $\vec u\,(-b, a)$ and
$\vec v\,(-b', a')$, the \omterm{def:g11:vect:det}{determinant} is $(-b)a' - (-b')a = ab' - a'b$.
If the lines are not parallel, solving the two \omterm{def:g10:algebra:equation}{equations} simultaneously
(by substitution or combination) leads to a unique solution: the system
behaves like a \omterm{def:g11:vect:det}{nonzero-determinant} $2 \times 2$ system.
\end{proof}

\begin{method}[Intersection of two lines]\label{met:g11:vect:intersection}
Solve the system of the two \omterm{def:g10:algebra:equation}{equations}: isolate one unknown in one \omterm{def:g10:algebra:equation}{equation}
and substitute into the other, or combine the \omterm{def:g10:algebra:equation}{equations} to eliminate one
unknown. Always check the resulting point in \emph{both} \omterm{def:g10:algebra:equation}{equations}.
\end{method}

\begin{example}\label{ex:g11:vect:intersection}
Intersect $d: x - 3y + 5 = 0$ and $d': 2x + y - 4 = 0$. From $d'$:
$y = 4 - 2x$. Substituting into $d$:
\[
x - 3(4 - 2x) + 5 = 0
\iff x - 12 + 6x + 5 = 0
\iff 7x = 7
\iff x = 1,
\]
then $y = 2$. The lines meet at $(1, 2)$. Check in $d$: $1 - 6 + 5 = 0$.
\end{example}

\section{Exercises}

\begin{exercise}[$\star$]\label{exo:g11:vect:1}
Given $A(2, -1)$, $B(5, 3)$ and $C(-1, 1)$, compute the \omterm{def:g10:coordgeom:system}{coordinates} of
$\vect{AB}$, $\vect{AC}$, $\vect{AB} + \vect{AC}$ and $2\vect{AB} -
3\vect{AC}$, and the \omterm{prop:g10:coordgeom:midpoint}{midpoint} of $[BC]$.
\end{exercise}

\begin{exercise}[$\star$]\label{exo:g11:vect:2}
Are the \omterm{def:g11:vect:vector}{vectors} \omterm{def:g11:vect:collinear}{collinear}?
\[
\vec u\,(4, -6) \text{ and } \vec v\,(-2, 3); \qquad
\vec u\,(2, 5) \text{ and } \vec v\,(3, 7); \qquad
\vec u\,(0, 3) \text{ and } \vec v\,(0, -8).
\]
\end{exercise}

\begin{exercise}[$\star$]\label{exo:g11:vect:3}
Find a cartesian \omterm{def:g10:algebra:equation}{equation} of the line through $A(2, 1)$ with \omterm{def:g11:vect:direction}{direction
vector} $\vec u\,(3, -1)$, then of the line through $B(0, 4)$ and
$C(2, 0)$.
\end{exercise}

\begin{exercise}[$\star$]\label{exo:g11:vect:4}
For the line $d: 3x - 2y + 6 = 0$, give a \omterm{def:g11:vect:direction}{direction vector}, the
\omterm{def:g10:numbers:interunion}{intersections} with the two coordinate axes, and decide whether the points
$P(2, 6)$ and $Q(1, 4)$ belong to $d$.
\end{exercise}

\begin{exercise}[$\star\star$]\label{exo:g11:vect:5}
Are the points $A(1, 2)$, $B(4, 8)$, $C(-2, -4)$ aligned? Same question
for $D(0, 1)$, $E(2, 4)$, $F(5, 9)$.
\end{exercise}

\begin{exercise}[$\star\star$]\label{exo:g11:vect:6}
Determine the \omterm{def:g10:numbers:interunion}{intersection} point, if any:
\[
d: 2x + y - 7 = 0 \ \text{ and }\ d': x - y + 1 = 0;
\qquad
e: x - 2y + 3 = 0 \ \text{ and }\ e': -2x + 4y + 1 = 0 .
\]
\end{exercise}

\begin{exercise}[$\star\star$]\label{exo:g11:vect:7}
Give a parametric representation of the line $d: 2x - y + 3 = 0$, and a
cartesian \omterm{def:g10:algebra:equation}{equation} of the line
$\bigl\{(1 + 2t,\ -3 + t) : t \in \R\bigr\}$.
\end{exercise}

\begin{exercise}[$\star\star$]\label{exo:g11:vect:8}
Find the \omterm{def:g10:algebra:equation}{equation} of the line through $P(1, -2)$ parallel to
$d: 4x - 3y + 2 = 0$, and the value of $m$ for which the line
$mx + 2y - 5 = 0$ is parallel to $d$.
\end{exercise}

\begin{exercise}[$\star\star$]\label{exo:g11:vect:9}
Let $A(-1, 0)$, $B(3, 2)$ and $C(1, -4)$. Find a cartesian \omterm{def:g10:algebra:equation}{equation} of the
\emph{\omterm{def:g10:stats:median}{median}} of the triangle $ABC$ issued from $C$ (the line joining $C$
to the \omterm{prop:g10:coordgeom:midpoint}{midpoint} of $[AB]$).
\end{exercise}

\begin{exercise}[$\star\star$]\label{exo:g11:vect:10}
For which value(s) of the parameter $m$ are the \omterm{def:g11:vect:vector}{vectors}
$\vec u\,(m, 2)$ and $\vec v\,(3, m - 1)$ \omterm{def:g11:vect:collinear}{collinear}?
\end{exercise}

\begin{exercise}[$\star\star\star$]\label{exo:g11:vect:11}
Let $ABCD$ be a parallelogram, $I$ the \omterm{prop:g10:coordgeom:midpoint}{midpoint} of $[AB]$, and $J$ the
point defined by $\vect{DJ} = \frac23 \vect{DI}$. Working in the
\omterm{def:g10:coordgeom:system}{coordinate system} where $A(0,0)$, $B(1,0)$, $D(0,1)$ (so $C(1,1)$):
\begin{enumerate}
  \item Compute the \omterm{def:g10:coordgeom:system}{coordinates} of $I$ and $J$.
  \item Show that $A$, $J$ and $C$ are aligned. (A nice fact: the line
        $(DI)$ cuts the diagonal $(AC)$ one third of the way up.)
\end{enumerate}
\end{exercise}

\section{Problem: Collision courses and clever coordinates}

\begin{problem}[{Weekend problem --- crossing paths is not
colliding: closest approach at sea, and theorems proved by
choosing the right frame}]
\label{pb:g11:vect:1}
Two ships' straight courses cross --- must the ships collide? Not
unless they reach the crossing \emph{at the same time}: paths are
sets of points, motions are parametric, and confusing the two has
sunk real ships. This problem runs a small vessel-traffic service
on parametric lines (\cref{def:g11:vect:parametric}), then turns
to pure geometry with the other superpower of \omterm{def:g10:coordgeom:system}{coordinates}:
\emph{choosing} them cleverly, as in \cref{exo:g11:vect:11}, to
make classical theorems compute themselves.

\medskip
\noindent\textbf{Part I --- Lines, three ways.}
\begin{enumerate}
  \item Give a parametric representation of the line through
        $A(1, 2)$ with \omterm{def:g11:vect:direction}{direction vector} $\vec u(3, -1)$, then
        eliminate the parameter to obtain a Cartesian \omterm{def:g10:algebra:equation}{equation}.
  \item For the line $2x - 5y + 3 = 0$: read off a \omterm{def:g11:vect:direction}{direction
        vector} (\cref{thm:g11:vect:cartesian}), find one point,
        and write a parametric representation.
  \item Compute the \omterm{def:g10:numbers:interunion}{intersection} point of the two lines of
        questions 1 and 2 (\cref{met:g11:vect:intersection}).
  \item Confirm with the \omterm{def:g11:vect:det}{determinant}
        (\cref{thm:g11:vect:collinearity}) that those two lines
        had to intersect; then show that $2x - 5y + 3 = 0$ and
        $-4x + 10y + 1 = 0$ are strictly parallel.
  \item Are $A(2, 1)$, $B(5, 3)$, $C(11, 7)$ aligned
        (\cref{met:g11:vect:alignment})?
\end{enumerate}

\medskip
\noindent\textbf{Part II --- The vessel-traffic service.}
At time $t$ (hours), ship $P$ is at $(2t,\ 3 + t)$ and ship $Q$
at $(10 - t,\ 2t - 1)$ (distances in nautical miles).
\begin{enumerate}[resume]
  \item Find the Cartesian \omterm{def:g10:algebra:equation}{equations} of the two \emph{paths},
        and compute where the paths cross.
  \item At what time does each ship pass the crossing point? Do
        the ships collide? State the general warning in one
        sentence: crossing paths against meeting motions.
  \item Compute the squared distance $D^2(t)$ between the ships
        at time $t$, and reduce it to a quadratic in $t$. At
        what time are the ships closest, and how close do they
        come (\cref{pb:g11:quad:1} vertex technology)?
  \item Regulations demand a separation of at least $1$ mile.
        Is it violated? If so, solve $D^2(t) < 1$ and give the
        duration of the violation.
  \item Explain why, for \emph{any} two ships on straight
        courses at constant speeds, $D^2(t)$ is always a
        \omterm{def:g11:quad:quadratic}{quadratic function} of $t$ --- so ``closest point of
        approach'' is always a vertex computation.
  \item Interception: a patrol boat starts at the origin at
        $t = 0$ with constant velocity $(a, b)$ and must meet
        ship $P$ exactly at $t = 3$. Compute the required
        $(a, b)$ and the patrol boat's speed.
\end{enumerate}

\medskip
\noindent\textbf{Part III --- Clever \omterm{def:g10:coordgeom:system}{coordinates}.}
Work in the frame of \cref{exo:g11:vect:11}: the parallelogram
$ABCD$ becomes $A(0,0)$, $B(1,0)$, $C(1,1)$, $D(0,1)$; let $I$
be the \omterm{prop:g10:coordgeom:midpoint}{midpoint} of $[AB]$ and $K$ the \omterm{prop:g10:coordgeom:midpoint}{midpoint} of $[CD]$.
\begin{enumerate}[resume]
  \item Compute a Cartesian \omterm{def:g10:algebra:equation}{equation} of the line $(DI)$.
  \item Intersect $(DI)$ with the diagonal $(AC)$: recover
        \cref{exo:g11:vect:11}'s fact --- the cut lies exactly
        one third of the way up the diagonal.
  \item Now the line $(BK)$: find its \omterm{def:g10:algebra:equation}{equation}, intersect with
        $(AC)$, and conclude the full classical theorem:
        \emph{the cevians $(DI)$ and $(BK)$ trisect the diagonal
        $(AC)$}.
  \item Why was it legitimate to prove the theorem in this one
        special frame? (What does the choice of \omterm{def:g10:coordgeom:system}{coordinates}
        preserve, and what does the statement involve?) Answer
        in two or three sentences.
  \item One more trisection for practice: let $E$ be the
        \omterm{prop:g10:coordgeom:midpoint}{midpoint} of $[BC]$. Where does the line $(AE)$ cut the
        \emph{other} diagonal $(DB)$?
\end{enumerate}

\medskip
\noindent\textbf{Part IV --- Position puzzles.}
\begin{enumerate}[resume]
  \item Are the three lines $x + y = 4$, \ $2x - y = 2$, \
        $x - 2y = -2$ concurrent? (Intersect two, test the
        third.)
  \item For each real $m$, let $d_m$ be the line
        $mx - y + 2 - 3m = 0$. Show that \emph{every} $d_m$
        passes through one \omterm{pb:g10:functions:1}{fixed point}. (Group the terms in
        $m$.)
  \item In the family $(d_m)$: which member is parallel to
        $y = 2x + 7$? Which passes through the origin?
  \item Finale --- the three costumes of a line: Cartesian
        (membership tests in one substitution), reduced (\omterm{def:g10:reffunc:affine}{slope}
        and \omterm{def:g10:functions:graph}{graph} at sight), parametric (motion, timing,
        interception); the \omterm{def:g11:vect:det}{determinant} as parallelism oracle;
        and the well-chosen frame as a theorem factory. One
        sentence each.
\end{enumerate}
\end{problem}
